\section{Literature Review}
\label{sec:lit-review}
\subsection{Governance Fragmentation Under the Paris Agreement}
\label{subsec:lit-governance}
International carbon markets have gone through boom, collapse, and post-Paris reconfiguration \cite{michaelowa-evolution}. The Paris Agreement shifted from top-down targets to bottom-up nationally determined contributions (NDCs). Article 6 creates two pathways. Article 6.2 allows bilateral transfers of "internationally transferred mitigation outcomes" (ITMOs). Article 6.4 establishes a centrally governed crediting mechanism \cite{ahonen-governance}.

Ahonen et al. \cite{ahonen-governance} call the current landscape "fragmented governance." Multiple public and private standards compete across market segments. This fragmentation hurts small-scale SSA projects most. Developers navigate competing standards (e.g., Gold Standard, Verra), divergent accounting rules, and uncertain demand. Host countries now bear "an unprecedented task of ensuring the environmental integrity and robust accounting of ITMOs." Many SSA governments lack this capacity.

\subsection{Environmental Integrity Risks}
\label{subsec:lit-risks}
Schneider and La Hoz Theuer \cite{schneider-integrity} define environmental integrity for carbon markets: transfers must lead to the same or lower global emissions. They identify four factors: robust accounting, unit quality, NDC ambition, and future mitigation incentives. Unit quality means one transferred unit represents at least 1 \tcooe of actual reduction. Quality matters most when emission sources fall outside NDC scope or when the NDC of a country requires no further mitigation. Both conditions apply to many SSA contexts.

There also exists a "hot air" risk: surplus units without real mitigation. La Hoz Theuer et al. \cite{theuer-transfers} compared NDC targets with business-as-usual emissions across 131 countries. Hot air volume equalled total pledged mitigation. In their low mitigation scenario, hot air reaches 5.4 G\tcooe. Between 36 and 52 percent of assessed NDCs could contain hot air.

Single-year NDC targets are common in SSA. They create accounting risks. Siemons and Schneider \cite{siemons-accounting} compared two approaches: averaging counts credits in the target year based on multi-year average transfers and multi-year approaches using emissions trajectories. Averaging is problematic. Year-on-year emission fluctuations magnify ITMO volumes needed, creating uncertainty for developers. Multi-year approaches are more robust but require technical capacity to establish credible trajectories.

\subsection{Proposed Solutions and Their Limits}
\label{subsec:lit-solutions}
Scholars have proposed technical fixes. Michaelowa et al. \cite{michaelowa-towards} introduced "ambition coefficients," declining multipliers applied to business-as-usual baselines, aligning crediting with net-zero pathways while respecting Common But Differentiated Responsibilities and Respective Capabilities (CBDR-RC). Developed countries would stop generating avoidance credits by 2035–2040, and developing countries could continue until 2050–2070. However, the approach needs regulatory coordination, technical capacity, and political consensus that most SSA contexts lack. The proposal also offers little guidance for small-scale commercial projects.

The World Bank has tried aggregation through carbon funds, alternating between "agenda setter" and "regime follower." Unfortunately, in the post-Paris period, its activities are "smaller, and...less transparent, than during the Kyoto Protocol period" \cite{michaelowa-catalysing}. This raises concerns about credit quality and replicability.

\subsection{Carbon Market Experience of Mini-Grids in Sub-Saharan Africa}
\label{subsec:vcm-experience}
The methodology of the Clean Development Mechanism (CDM), AMS-III.BB, "Electrification of communities through grid extension or construction of new mini-grids," was adopted in 2012 and has since been updated twice \cite{unfccc-ams-iii-bb}. Version 3.0 requires continuous metering and user surveys, caps annual emission reductions at 60,000 \tcooe, and requires at least 75 percent of users must be households. This methodology provided the technical foundation for later mini-grid carbon projects under Gold Standard, Verra, and the Renewvia Environmental Exchange, including the atmosfair-MiniGridCo PoAs in Nigeria.

Michaelowa et al. \cite{michaelowa-mobilising} provide the most comprehensive assessment of SSA carbon markets, identifying 457 private-sector mitigation projects. They note that SSA was "essentially bypassed by the CDM until the late 2000s," but PoAs increased participation. The SSA share in PoAs was ten times higher than for single projects. More recently however, implementation has stalled. Only 22 percent of single CDM projects and 16 percent of PoA component activities have issued credits. "Some may have paused or stopped operation due to the low market price of Certified Emission Reductions (CERs)."

Revenue leakage makes the problem worse. A Power Africa white paper found that only 30 percent of carbon credit revenues reach project developers, with the rest going to intermediaries, certification fees, and transaction costs \cite{power-africa-credits}. More recently, SSA governments have demanded an additional cut of the proceeds. For example, Kenya's Climate Change Act added a 25 percent tax on carbon projects (40 percent for land-based projects) \cite{kenya-climate-change-act}.

Ethiopian PoAs show both potential and limits. They enabled better CDM access but remain "restricted to sustainable energy access in rural areas due to the prohibitively low emissions baseline for grid electricity" \cite{michaelowa-mobilising}. South Africa's Renewable Energy Independent Power Producer Procurement Programme (REIPPPP) succeeded but is not replicable: "South Africa is particularly well positioned...As other SSA countries often still lack these capacities, REIPPPP cannot be easily replicated" \cite{michaelowa-mobilising}.

Kenya illustrates the gap between theoretical potential and practical implementation. The country has a Climate Change Act (2016, amended 2023), Carbon Markets Regulations (2024), and a National Carbon Registry. Yet it has approved zero Article 6 projects \cite{kenya-zero-approvals}. The government fixed Corresponding Adjustment prices at USD 4.00 per \tcooe and strictly limits Letters of Authorisation to preserve its "carbon budget" \cite{kenya-carbon-budget}.

This gap led to the infamous collapse of KOKO Networks in February 2026. The company served 1.5 million Kenyan households with bioethanol cooking fuel, depended on carbon revenues to subsidize fuel prices. The government refused a Letter of Authorisation, barring monetization under Article 6.2 or the Carbon Offsetting and Reduction Scheme for International Aviation (CORSIA). KOKO entered administration. Seven hundred employees lost jobs, and 1.5 million households returned to charcoal and kerosene (\cite{koko-disrupt,koko-econews}. One analyst said this "highlights the extreme fragility of business models that rely on the complex and often shifting regulatory landscapes of international carbon markets" \cite{koko-econews}. The case validates Schneider and La Hoz Theuer's \cite{schneider-integrity} theoretical risk. Even real, additional, verifiable reductions cannot succeed without host government authorization in the Article 6.2 markets.

\subsection{Research Gap and Contribution}
\label{subsec:lit-gap}
The literature establishes three insights. First, carbon markets face integrity risks: hot air, additionality failure, accounting complexity. These threaten to increase global emissions \cite{schneider-integrity,theuer-transfers,siemons-accounting}. Second, proposed governance solutions assume institutional capacities most SSA contexts lack \cite{michaelowa-towards,michaelowa-catalysing}. Third, evidence shows that while PoAs increased SSA participation, most projects remain stalled. Low prices, revenue leakage, and host government authorization barriers all play a role \cite{michaelowa-mobilising,power-africa-credits,koko-disrupt}.

No empirical study has systematically examined how these risks and solutions manifest in rural solar mini-grids. Current conditions include 1–5 USD per credit, fragmented governance, and uncertain authorization. This paper asks: does carbon revenue influence mini-grid investment decisions? How do transaction costs compare to revenues at scales below 200 kW? What conditions enable successful integration at scale, including host government authorization?